\documentclass[12pt,a4paper]{article}
\newcommand{\LabNumber}{6}
\newcommand{\LabTitle}{Funcții hash și semnături digitale:\\ RSA și ElGamal cu HAVAL-224,4}
\newcommand{\LabTitleShort}{Lucrarea de laborator nr. 6 — Criptografie \& Securitate Informațională}
% ---------------------------------------------------------------------------
% Common preamble for the CS laboratory reports (compile with XeLaTeX)
% Strunga Daniel-Ioan, FAF-242 — https://github.com/don09dandidi/CS
% ---------------------------------------------------------------------------
\usepackage{fontspec}
\setmainfont{Latin Modern Roman}
\setsansfont{Latin Modern Sans}
\setmonofont{DejaVu Sans Mono}[Scale=0.85]
\usepackage[english]{babel}
\babelprovide[import, main]{romanian}
\usepackage[a4paper, margin=2.5cm, headheight=15pt]{geometry}
\usepackage{amsmath, amssymb}
\usepackage{booktabs, array, longtable, tabularx}
\usepackage{xcolor}
\usepackage{listings}
\usepackage{fvextra}
\usepackage[most]{tcolorbox}
\usepackage{fancyhdr}
\usepackage{tikz}
\usetikzlibrary{arrows.meta, positioning, shapes.geometric, fit, calc}
\usepackage{enumitem}
\usepackage{caption}
\usepackage{float}
\usepackage[hidelinks]{hyperref}
\usepackage{url}

\definecolor{codebg}{RGB}{248,248,250}
\definecolor{codeframe}{RGB}{200,200,210}
\definecolor{kw}{RGB}{0,0,160}
\definecolor{cmt}{RGB}{0,120,0}
\definecolor{str}{RGB}{160,30,30}
\definecolor{termbg}{RGB}{30,32,40}
\definecolor{termfg}{RGB}{225,228,235}
\definecolor{termbar}{RGB}{60,63,75}
\definecolor{boxblue}{RGB}{230,238,250}

\lstdefinestyle{py}{
  language=Python, basicstyle=\ttfamily\small, backgroundcolor=\color{codebg},
  frame=single, rulecolor=\color{codeframe}, numbers=left, numberstyle=\tiny\color{gray},
  keywordstyle=\color{kw}\bfseries, commentstyle=\color{cmt}\itshape, stringstyle=\color{str},
  breaklines=true, showstringspaces=false, columns=fullflexible, keepspaces=true, tabsize=4,
  extendedchars=true, inputencoding=utf8
}
\lstdefinestyle{sh}{
  language=bash, basicstyle=\ttfamily\small, backgroundcolor=\color{codebg},
  frame=single, rulecolor=\color{codeframe}, numbers=left, numberstyle=\tiny\color{gray},
  keywordstyle=\color{kw}\bfseries, commentstyle=\color{cmt}\itshape, stringstyle=\color{str},
  breaklines=true, showstringspaces=false, columns=fullflexible, keepspaces=true
}
\lstdefinestyle{wl}{
  basicstyle=\ttfamily\small, backgroundcolor=\color{codebg},
  frame=single, rulecolor=\color{codeframe}, numbers=left, numberstyle=\tiny\color{gray},
  morekeywords={RandomPrime,RandomInteger,PowerMod,Mod,GCD,FromDigits,IntegerDigits,
    ToCharacterCode,FromCharacterCode,IntegerLength,BitLength,Module,Print,Hash,Length},
  keywordstyle=\color{kw}\bfseries, commentstyle=\color{cmt}\itshape,
  morecomment=[s]{(*}{*)}, breaklines=true, columns=fullflexible, keepspaces=true
}

% Console session box (as in the Lab 2 report)
\newtcolorbox{terminal}[1][Sesiune Terminal Consolă]{
  enhanced, breakable, colback=termbg, colframe=termbar, coltext=termfg,
  fonttitle=\small\ttfamily\bfseries, title={\textbullet\textbullet\textbullet\quad #1},
  left=4pt, right=4pt, top=2pt, bottom=2pt, boxrule=0.6pt, arc=2pt
}
\DefineVerbatimEnvironment{termout}{Verbatim}{fontsize=\scriptsize, breaklines=true, breakanywhere=true}

\newtcolorbox{infobox}[1]{
  enhanced, breakable, colback=boxblue, colframe=kw!60, fonttitle=\bfseries,
  title={#1}, boxrule=0.6pt, arc=2pt
}

\newcommand{\bignum}[1]{\par\nobreak\vspace{1pt}{\noindent\ttfamily\footnotesize\seqsplit{#1}\par}\vspace{5pt}}
\usepackage{seqsplit}

\pagestyle{fancy}
\fancyhf{}
\renewcommand{\headrulewidth}{0.4pt}
\fancyhead[L]{\footnotesize\textit{\LabTitleShort}}
\fancyhead[R]{\footnotesize strunga.daniel-ioan}
\fancyfoot[C]{\thepage}

\newcommand{\makereporttitle}{%
\begin{titlepage}
  \centering
  {\large\bfseries UNIVERSITATEA TEHNICĂ A MOLDOVEI\par}
  \vspace{2pt}
  {\normalsize Facultatea Calculatoare, Informatică și Microelectronică\par}
  {\normalsize Departamentul Ingineria Software și Automatică\par}
  \vspace{5cm}
  {\LARGE\bfseries LUCRAREA DE LABORATOR NR. \LabNumber\par}
  \vspace{8pt}
  {\Large\bfseries \LabTitle\par}
  \vspace{10pt}
  {\normalsize Disciplina: Criptografie și Securitate Informațională\par}
  \vfill
  \begin{tabular*}{\textwidth}{@{}l@{\extracolsep{\fill}}r@{}}
    \textbf{Autor:} & \textbf{Verificat de:} \\
    Daniel-Ioan Strunga & Lector / Supervizor \\
    Student gr. FAF-242 & Maia Zaica \\
  \end{tabular*}
  \vspace{3cm}

  {\normalsize Chișinău — 2026\par}
\end{titlepage}
}


\begin{document}
\makereporttitle
\tableofcontents
\newpage

\section{Introduction and Laboratory Objectives}
A \textbf{digital signature} gives a document the properties of a handwritten signature —
authenticity, integrity and non-repudiation — using public-key cryptography. Because public-key
operations are slow and work on numbers smaller than the modulus, the document itself is not
signed: a \textbf{cryptographic hash function} first compresses it into a short fixed-length
digest, and the digest is signed (\emph{hash-then-sign}).

The variant is determined by the student's number in the group list, $k = <<k>>$:
\[
  i = (k \bmod 24) + 1 = (<<k>> \bmod 24) + 1 = <<hash_index>>
\]
In \emph{both} lists of the assignment (Task 2 and Task 3) index 23 corresponds to
\textbf{HAVAL-224,4} — the HAVAL hash with a 224-bit digest and 4 passes.

Objectives:
\begin{enumerate}
  \item \textbf{Task 2} — generate RSA keys with $n \ge 3072$ bits, sign the message $m$ obtained
        in Laboratory Work No.~2 and validate the signature, using HAVAL-224,4.
  \item \textbf{Task 3} — sign and validate the same message with the \textbf{ElGamal} signature
        scheme, using the given 2048-bit prime $p$ and generator $g=2$, and HAVAL-224,4.
  \item Use the decimal representation of the hash value and comment every step of the algorithms.
\end{enumerate}

\begin{infobox}{Tooling note}
The guide suggests Wolfram Alpha / Mathematica. Wolfram's \texttt{Hash[]} function does not
implement HAVAL, so the digest was computed with a pure-Python implementation of HAVAL
(\texttt{lab6/haval.py}, checked against PHP's \texttt{hash("haval224,4", ...)}) and then
converted from hexadecimal to decimal. The modular arithmetic was performed with
\texttt{lab6/digital\_signatures.py} (Python integers, the same operations as \texttt{RandomPrime},
\texttt{RandomInteger}, \texttt{PowerMod}, \texttt{Mod}); the equivalent Wolfram Language code is
given in Appendix~\ref{app:wl}.
\end{infobox}

\section{The Message}
The message $m$ is the plaintext recovered in Laboratory Work No.~2 (Variant 22, a fragment of
D.~Kahn's \emph{The Codebreakers} about William F.~Friedman). It is stored in
\texttt{lab6/message\_lab2.txt} as UTF-8 text, three paragraphs separated by a newline, without a
trailing newline: \textbf{<<msg_len>> bytes}.
\begin{infobox}{$m$ (beginning and end)}
\small\itshape While Herbert Yardley may be the best known cryptologist, uncontestably the greatest
is William Frederick Friedman. Unlike his contemporary, his eminence is due most emphatically to
what he did. [\dots] A cardinal article of faith with him was that a well-executed sales campaign
could put across almost anything.
\end{infobox}

\section{The HAVAL Hash Function}
\subsection{Description}
HAVAL was proposed by Y.~Zheng, J.~Pieprzyk and J.~Seberry (AUSCRYPT'92). It is a
Merkle–Damgård hash of the MD family, with two parameters: the number of passes (3, 4 or 5) and
the digest length (128, 160, 192, 224 or 256 bits).
\begin{itemize}
  \item \textbf{State:} eight 32-bit words, initialised with the first 256 bits of the fractional
        part of $\pi$:\\ {\small\texttt{243F6A88 85A308D3 13198A2E 03707344 A4093822 299F31D0 082EFA98 EC4E6C89}}.
  \item \textbf{Padding:} append byte \texttt{0x01}, then zeros up to 944 bits mod 1024, then
        16 bits encoding VERSION=1, PASS and FPTLEN, then the 64-bit message length
        (little-endian). The message is processed in \textbf{1024-bit blocks} of 32 words.
  \item \textbf{Compression:} each pass has 32 steps. Step $j$ of pass $r$ updates one register:
        \[ T = F_r\big(\phi_{P,r}(x_6,\dots,x_0)\big),\qquad
           x_7 \leftarrow (T \ggg 7) + (x_7 \ggg 11) + W_{\mathrm{ord}_r(j)} + K_{r,j}, \]
        where $F_1\dots F_5$ are highly non-linear Boolean functions of 7 variables, $\phi_{P,r}$
        is a fixed permutation of the inputs (depending on the number of passes $P$),
        $\mathrm{ord}_r$ a word order and $K_{r,j}$ further words of $\pi$ (none in pass 1). After
        all passes the result is added word-wise to the old state.
  \item \textbf{Tailoring (224 bits):} the 256-bit state is folded to 224 bits by adding 5- and
        4-bit fields of the last word $H_7$ to $H_0\dots H_6$.
\end{itemize}
HAVAL-224,4 therefore has $4\times32 = 128$ steps per block. Its security is lower than that of
SHA-2/SHA-3 (collisions for 3-pass HAVAL were found in 2004), so in practice it should not be
used for new signatures; here it is used as required by the variant.

\subsection{Implementation and Validation}
\begin{lstlisting}[style=py, caption={Compression function of HAVAL (\texttt{haval.py}, excerpt)}]
def _compress(state, block: bytes, passes: int):
    w = [int.from_bytes(block[4*i:4*i+4], "little") for i in range(32)]
    t = list(state)
    for p in range(passes):
        f, phi = _F[p], _PHI[passes][p]
        for s in range(32):
            x = [t[(j - s) % 8] for j in range(8)]        # x7..x0 rotate each step
            temp = f(*[x[j] for j in phi]) & MASK
            new = (_rotr(temp, 7) + _rotr(x[7], 11) + w[ORDER[p][s]]) & MASK
            if p > 0:
                new = (new + CONST[p][s]) & MASK
            t[(7 - s) % 8] = new
    return [(a + b) & MASK for a, b in zip(state, t)]
\end{lstlisting}
The implementation was compared with PHP's \texttt{hash()} for HAVAL-224 and HAVAL-256 with 3,
4 and 5 passes on 10 inputs of different lengths (including the block boundaries 117–129 bytes):
all 60 digests are identical. For example
\texttt{HAVAL-224,4("")} $=$ \texttt{3e56243275b3b81561750550e36fcd676ad2f5dd9e15f2e89e6ed78e}.

\subsection{Digest of the Message}
\noindent\textbf{Hexadecimal digest} $\mathrm{HAVAL\text{-}224,4}(m)$:
\bignum{<<h_hex>>}
\noindent\textbf{Decimal value} $h$ (<<h_bits>> bits), used in both signatures:
\bignum{<<h>>}
\textbf{Avalanche effect:} replacing a single letter (the first ``Friedman'' $\to$ ``Fridman'')
gives the completely different digest
\bignum{<<hf_hex>>}

\section{Task 2 — RSA Digital Signature}
\subsection{Algorithm}
\begin{enumerate}
  \item \textbf{Key generation:} random primes $p, q$ of 1536 bits each, $n = pq$ (3072 bits),
        $\varphi(n) = (p-1)(q-1)$, $e = 65537$ with $\gcd(e,\varphi(n)) = 1$,
        $d = e^{-1} \bmod \varphi(n)$. Public key $(e,n)$, private key $(d,n)$.
  \item \textbf{Signing} (with the private key): $h = H(m)$, \quad $s = h^{d} \bmod n$.
        The signed message is the pair $(m, s)$.
  \item \textbf{Verification} (with the public key): compute $h = H(m)$ from the received message
        and $h' = s^{e} \bmod n$; the signature is valid iff $h' = h$.
\end{enumerate}
Since $h < 2^{224} < n$, the digest is signed directly as a number (\emph{textbook} RSA signature;
real systems apply PKCS\#1 v1.5 or PSS padding first).

\subsection{Key Generation}
\noindent\textbf{Prime} $p$ (1536 bits):
\bignum{<<rsa.p>>}
\noindent\textbf{Prime} $q$ (1536 bits):
\bignum{<<rsa.q>>}
\noindent\textbf{Modulus} $n = p\cdot q$ (\textbf{<<rsa.n_bits>> bits}):
\bignum{<<rsa.n>>}
\noindent\textbf{Euler's function} $\varphi(n)$:
\bignum{<<rsa.phi>>}
\noindent\textbf{Public exponent} $e = <<rsa.e>>$; \textbf{private exponent} $d = e^{-1} \bmod \varphi(n)$:
\bignum{<<rsa.d>>}

\subsection{Signing and Verification}
\noindent\textbf{Signature} $s = h^{d} \bmod n$:
\bignum{<<rsa.s>>}
\noindent\textbf{Verification} $h' = s^{e} \bmod n$:
\bignum{<<rsa.h_check>>}
$h' = h$ $\Rightarrow$ the signature is \textbf{valid} (<<rsa.valid>>). For the modified message
$m'$, $h' \ne H(m')$ $\Rightarrow$ the signature is rejected (check $=$ <<rsa.valid_fake>>).

\section{Task 3 — ElGamal Digital Signature}
\subsection{Algorithm}
Public parameters: the given 2048-bit prime $p$ and the generator $g = 2$.
\begin{enumerate}
  \item \textbf{Key generation:} private key $x$ random, $1 < x < p-1$; public key
        $y = g^{x} \bmod p$.
  \item \textbf{Signing:} choose a random $k$, $1<k<p-1$, with $\gcd(k, p-1) = 1$ (so that
        $k^{-1} \bmod (p-1)$ exists) and compute
        \[ r = g^{k} \bmod p, \qquad s = (h - x\,r)\,k^{-1} \bmod (p-1). \]
        The signature is the pair $(r,s)$.
  \item \textbf{Verification:} accept iff $0<r<p$ and
        \[ g^{h} \equiv y^{r}\, r^{s} \pmod p. \]
\end{enumerate}
\textbf{Correctness:} from the definition of $s$, $h \equiv x r + k s \pmod{p-1}$, therefore
$g^{h} = g^{xr}\,g^{ks} = (g^{x})^{r}(g^{k})^{s} = y^{r} r^{s} \pmod p$ (Fermat's little theorem
allows exponents to be reduced modulo $p-1$). The value $k$ must be secret and unique for every
signature — reusing it reveals the private key $x$.

\subsection{Values}
\noindent\textbf{Prime} $p$ (given, 2048 bits), $g = 2$:
\bignum{<<P>>}
\noindent\textbf{Private key} $x$:
\bignum{<<eg.x>>}
\noindent\textbf{Public key} $y = g^{x} \bmod p$:
\bignum{<<eg.y>>}
\noindent\textbf{Random} $k$ ($\gcd(k,p-1) = <<eg.gcd_k>>$):
\bignum{<<eg.k>>}
\noindent$k^{-1} \bmod (p-1)$:
\bignum{<<eg.k_inv>>}
\noindent\textbf{Signature} $r = g^{k} \bmod p$:
\bignum{<<eg.r>>}
\noindent\textbf{Signature} $s = (h - x r)\,k^{-1} \bmod (p-1)$:
\bignum{<<eg.s>>}

\subsection{Verification}
\noindent Left side $g^{h} \bmod p$:
\bignum{<<eg.left>>}
\noindent Right side $y^{r} r^{s} \bmod p$:
\bignum{<<eg.right>>}
The two sides are equal $\Rightarrow$ the signature is \textbf{valid} (<<eg.valid>>). For the
modified message $m'$, $g^{H(m')} \ne y^{r} r^{s}$ and the signature is rejected
(check $=$ <<eg.valid_fake>>).

\section{Source Code and Repository}
\begin{center}
\url{https://github.com/don09dandidi/CS} \quad (directory \texttt{lab6/})
\end{center}
\begin{lstlisting}[style=py, caption={Signing and verification (excerpt from \texttt{digital\_signatures.py})}]
h = int(haval224_4(message_bytes), 16)          # decimal value of the digest

# RSA (n = 3072 bits)
s = pow(h, d, n)                                # sign
valid = pow(s, e, n) == h                       # verify

# ElGamal (p = 2048 bits, g = 2)
r = pow(G, k, P)
s = ((h - x * r) * pow(k, -1, P - 1)) % (P - 1)  # sign: (r, s)
valid = pow(G, h, P) == (pow(y, r, P) * pow(r, s, P)) % P
\end{lstlisting}

\subsection{Build and Execution Instructions}
\begin{lstlisting}[style=sh]
git clone https://github.com/don09dandidi/CS.git
cd CS/lab6
python3 haval.py < message_lab2.txt        # HAVAL-224,4 digest of the message
python3 digital_signatures.py --save       # keys, signatures, verification -> results.json
\end{lstlisting}

\section{Conclusion}
\begin{enumerate}
  \item For $k = 22$ the hash function is HAVAL-224,4; it was implemented from its specification,
        validated against an independent implementation, and produced the 224-bit digest of the
        Lab~2 message, used in decimal form.
  \item \textbf{RSA signature:} with a 3072-bit modulus the digest was signed with the private
        exponent and verified with the public one ($s^{e} \bmod n = h$).
  \item \textbf{ElGamal signature:} with the given 2048-bit prime and $g=2$ the signature
        $(r,s)$ was generated with a random $k$ coprime to $p-1$ and verified through
        $g^{h} \equiv y^{r} r^{s} \pmod p$.
  \item In both schemes a single-letter change of the message produces a completely different
        digest and the signature is rejected, demonstrating integrity protection. RSA signatures
        are deterministic and short to verify ($e$ small); ElGamal signatures are randomised and
        twice as long, and their security depends critically on never reusing $k$.
\end{enumerate}

\appendix
\section{Wolfram Language Code}\label{app:wl}
\lstinputlisting[style=wl]{wolfram/lab6.wl}

\begin{thebibliography}{9}
\bibitem{haval} Y.\ Zheng, J.\ Pieprzyk, J.\ Seberry, \emph{HAVAL — A One-Way Hashing Algorithm
  with Variable Length of Output}, AUSCRYPT'92, LNCS 718, 1993.
\bibitem{rsa} R.\ Rivest, A.\ Shamir, L.\ Adleman, \emph{A Method for Obtaining Digital Signatures
  and Public-Key Cryptosystems}, Communications of the ACM, 21(2), 1978.
\bibitem{elg} T.\ ElGamal, \emph{A Public Key Cryptosystem and a Signature Scheme Based on Discrete
  Logarithms}, IEEE Trans.\ Inf.\ Theory, 31(4), 1985.
\bibitem{kahn} David Kahn, \emph{The Codebreakers: The Story of Secret Writing}, Scribner, 1996.
\bibitem{stallings} William Stallings, \emph{Cryptography and Network Security: Principles and
  Practice}, 7th Edition, Pearson, 2017.
\end{thebibliography}
\end{document}
